\chapter{含参积分及广义含参积分}
\begin{definition}{含参积分}{}
    含参积分
    \[
        I(y)=\int_a^bf(x,y)\d x,\quad y\in[c,d].
    \]
    广义含参积分包括无穷积分和瑕积分，可写为
    \[
        I(y)=\int_a^\omega f(x,y)\d x:=\lim_{A\to\omega^-}\int_a^Af(x,y)\d x.
    \]
    无穷积分取$\omega=+\infty$；瑕积分可取有限端点$\omega$。下文假定$f$在每个紧矩形$[a,A]\times[c,d]$上连续，奇异性仅位于上端点。若有多个瑕点或两个无穷端点，须在内点处分段，分别检查收敛条件。
\end{definition}
\begin{definition}{一致连续}{}
    函数$f:\Omega\to\RR$一致连续，是指$\forall\varepsilon>0$，存在与点的位置无关的$\delta>0$，使所有$X,Y\in\Omega$只要$\norm{X-Y}<\delta$，便有$|f(X)-f(Y)|<\varepsilon$。
    \tcblower
    \paragraph{否定形式：}存在$\varepsilon_0>0$及两点列$X_k,Y_k\in\Omega$，满足$\norm{X_k-Y_k}\to0$，但对所有$k$有$|f(X_k)-f(Y_k)|\geq\varepsilon_0$。
\end{definition}
\begin{theorem}{一致连续的判定}{}
    定义在$\RR^n$的闭有界集上的连续函数一致连续。
\end{theorem}
\begin{definition}{一致收敛}{}
    设$I(y)$对每个$y\in[c,d]$均存在。若$\forall\varepsilon>0$，存在$N\in(a,\omega)$，使所有$N<A<\omega$及所有$y\in[c,d]$均满足
    \[
        \abs{\int_a^Af(x,y)\d x-I(y)}<\varepsilon,
    \]
    则称广义积分关于$y$一致收敛。这一定义同时适用于$\omega<\infty$和$\omega=+\infty$。
\end{definition}
\begin{theorem}{一致收敛的Weierstrass判别法}{}
    若对所有$x\in[a,\omega)$及$y\in[c,d]$都有$|f(x,y)|\leq F(x)$，其中$F\geq0$且$\int_a^\omega F(x)\d x<\infty$，则$\int_a^\omega f(x,y)\d x$关于$y$一致收敛。
\end{theorem}
\begin{theorem}{一致收敛的Dirichlet判别法}{}
    设$f,g$局部连续，并且
    \begin{compactenum}
        \item 存在常数$M$，使所有$A\in(a,\omega)$、$y\in[c,d]$满足$\abs{\int_a^Af(x,y)\d x}\leq M$；
        \item 对每个$y$，$g(x,y)$关于$x$单调，且$g(x,y)\to0$在$x\to\omega^-$时对$y$一致成立。
    \end{compactenum}
    则$\int_a^\omega f(x,y)g(x,y)\d x$关于$y$一致收敛。仅有逐个参数下的有界性或趋零性是不够的。
\end{theorem}
\begin{theorem}{一致收敛的Abel判别法}{}
    设$f,g$局部连续，$\int_a^\omega f(x,y)\d x$关于$y$一致收敛；对每个$y$，$g(x,y)$关于$x$单调，并且存在常数$M$使所有$x,y$满足$|g(x,y)|\leq M$。则$\int_a^\omega f(x,y)g(x,y)\d x$关于$y$一致收敛。
\end{theorem}
\section{积分号的可交换性}
交换极限、求导与积分必须逐项验证条件；另可参阅 Keith Conrad 的\href{https://kconrad.math.uconn.edu/blurbs/analysis/diffunderint.pdf}{Differentiating under the Integral Sign}，其中也讨论了条件收敛的情形。
\begin{theorem}{连续性}{}
    若$f$在$[a,b]\times[c,d]$上连续，则$I(y)=\int_a^bf(x,y)\d x$在$[c,d]$上连续，即
    \[
        \lim_{y\to y_0}\int_a^bf(x,y)\d x=\int_a^b\lim_{y\to y_0}f(x,y)\d x.
    \]
    \tcblower
    对广义积分，若$f$在每个截断矩形上连续，且$\int_a^\omega f(x,y)\d x$关于$y\in[c,d]$一致收敛，则$I(y)$连续，仍可交换参数极限与积分。
\end{theorem}
\begin{theorem}{可微性}{}
    若$f,f_y$在$[a,b]\times[c,d]$上连续，则$I$在$(c,d)$内可微，且
    \[
        I'(y)=\dd y\int_a^bf(x,y)\d x=\int_a^bf_y(x,y)\d x.
    \]
    \tcblower
    对广义积分，假设$f,f_y$在每个截断矩形上连续，$\int_a^\omega f(x,y_0)\d x$在至少一个$y_0\in[c,d]$收敛，且$\int_a^\omega f_y(x,y)\d x$关于$y\in[c,d]$一致收敛。则$I$存在并在$(c,d)$内可微，而且
    \[
        I'(y)=\int_a^\omega f_y(x,y)\d x.
    \]
    若只讨论某个内点附近，可在该点周围的紧参数区间上验证这些条件。
\end{theorem}
若有限积分上下限$\alpha,\beta\in C^1$，且$f,f_y$在相应积分区域的邻域内连续，则Leibniz公式为
\[
    \dd y\int_{\alpha(y)}^{\beta(y)}f(x,y)\d x
    =\int_{\alpha(y)}^{\beta(y)}f_y(x,y)\d x+f(\beta(y),y)\beta'(y)-f(\alpha(y),y)\alpha'(y).
\]
\begin{theorem}{积分性}{}
    若$f$在$[a,b]\times[c,d]$上连续，则
    \[
        \int_c^dI(y)\d y=\int_c^d\int_a^bf(x,y)\d x\nd y
        =\int_a^b\int_c^df(x,y)\d y\nd x.
    \]
    \tcblower
    若$f$在每个截断矩形上连续，且$\int_a^\omega f(x,y)\d x$关于$y\in[c,d]$一致收敛，则
    \[
        \int_c^dI(y)\d y=\int_c^d\int_a^\omega f(x,y)\d x\nd y
        =\int_a^\omega\int_c^df(x,y)\d y\nd x.
    \]
\end{theorem}
对于两个广义积分，一组便于使用的充分条件是：
\begin{compactenum}
    \item $f(x,y)$在$[a,+\infty)\times[c,+\infty)$上连续；
    \item 两个内层广义积分分别对另一参数的每个紧区间一致收敛；
    \item 下列两个非负累次积分中至少一个收敛到有限值：
    \[
        \int_c^{+\infty}\int_a^{+\infty}|f(x,y)|\d x\nd y,\qquad
        \int_a^{+\infty}\int_c^{+\infty}|f(x,y)|\d y\nd x.
    \]
\end{compactenum}
此时可交换积分次序：
\[
    \int_c^{+\infty}\int_a^{+\infty}f(x,y)\d x\nd y
    =\int_a^{+\infty}\int_c^{+\infty}f(x,y)\d y\nd x.
\]
绝对可积是控制换序的关键充分条件；仅有两个条件收敛的累次积分存在，不能保证它们相等。

Gamma与Beta积分的参数范围及关系可核对\href{https://dlmf.nist.gov/5.2}{NIST DLMF §5.2}和\href{https://dlmf.nist.gov/5.12}{§5.12}。
\begin{example}{Gamma函数}{}
	\begin{center}
		$\Gamma(x):=\int\zti t^{x-1}\e^{-t}\d t,\quad x>0.$
	\end{center}
	对$x>0$，分部积分中的两个端点项均为零，得到递推公式
	\begin{align*}
		\Gamma(1)   & =\int\zti\e^{-t}\d t=1.                                              \\
		\Gamma(x+1) & =\int\zti t^x\cdot\e^{-t}\d t                                        \\
					& =\eval{-t^x\e^{-t}}\zti+\int\zti\e^{-t}\cdot xt^{x-1}\d t=x\Gamma(x).
	\end{align*}
	对于正整数$n\geq1$，$\Gamma(n)=(n-1)!$。上述积分定义不适用于$x\leq0$，不能将其与解析延拓混为一谈。
	% $\Gamma(x+1)=x\Gamma(x),\Gamma(1)=1$，对正整数$n$,
	% \[
	% 	\Gamma(n)=\int_0^{+\infty}t^{n-1}\e^{-t}\d t=(n-1)!
	% \]

	\textbf{余元公式}
	\[
		\Gamma(x)\Gamma(1-x)=\frac{\pi}{\sin\pi x},\quad x\in(0,1).
	\]
	证明略，有$\Gamma(1/2)=\sqrt\pi$
\end{example}
\begin{example}{Beta函数}{}
	\begin{center}
		$\mathrm B(p,q):=\int_0^1t^{p-1}(1-t)^{q-1}\d t,\quad p,q>0.$ % =2\int_0^{\pi/2}\cos^{2p-1}\theta\sin^{2q-1}\theta\d\theta.$
	\end{center}
	在$p,q>0$条件下，与Gamma函数的关系为$\mathrm B(p,q)=\frac{\Gamma(p)\Gamma(q)}{\Gamma(p+q)}.$
\end{example}
\begin{example}{Poisson积分}{}
	\begin{center}
		$\int\zti\e^{-t^2}\d t=\frac{\sqrt\pi}2.$
	\end{center}
	\textbf{证明：}考虑积分的平方
	\begin{align*}
		\int\zti\e^{-x^2}\d x\int\zti\e^{-y^2}\d y=\int\zti\int\zti\e^{-x^2-y^2}\d x\nd y \\
		=\int_0^{\pi/2}\int\zti\e^{-r^2}\cdot r\d r\nd\theta=\frac\pi{2}\cdot\frac12\int\zti\e^{-s}\d s=\frac\pi{4}.
	\end{align*}
	\iffalse
		由$\ex>0,\forall x\in\RR$
		\[
			\iint_A\e^{-x^2-y^2}\d x\nd y\leqslant\iint_B\e^{-x^2-y^2}\d x\nd y\leqslant\iint_C\e^{-x^2-y^2}\d x\nd y
		\]
		其中
		\begin{align*}
			B & =\{(x,y)~|~0\leqslant x\leqslant R,0\leqslant y\leqslant R\}, \\
			A & =\{(x,y)~|~x^2+y^2\leqslant R^2,x\geqslant 0,y\geqslant 0\},  \\
			C & =\{(x,y)~|~x^2+y^2\leqslant 2R^2,x\geqslant 0,y\geqslant 0\}.
		\end{align*}
		又
		\begin{align*}
			\iint_A\e^{-x^2-y^2}\d x\nd y & =\int_0^{\pi/2}\int_0^R\e^{-\rho^2}\rho\d\rho\d\varphi=\frac{\pi}4(1-\e^{-R^2}), \\
			\iint_B\e^{-x^2-y^2}\d x\nd y & =\int_0^R\int_0^R\e^{-x^2}\e^{-y^2}\d x\nd y=\biggp{\int_0^R\e^{-t^2}\d t}^2,          \\
			\iint_C\e^{-x^2-y^2}\d x\nd y & =\frac{\pi}4(1-\e^{-2R^2}).
		\end{align*}
		当$R\to+\infty$，即得$\int\zti\e^{-t^2}\d t=\frac{\sqrt\pi}2$.
	\fi
\end{example}
\begin{example}{Dirichlet积分}{}
	\begin{center}
		$\int\zti\frac{\sin x}x\d x=\frac{\pi}2.$
	\end{center}
	\textbf{证明：}令
	\[
		F(y)=\int\zti\frac{\sin x}x\e^{-xy}\d x,\quad y\geqslant 0.
	\]
	被积函数在$x=0$可连续延拓，但这本身不足以推出广义积分$F(y)$连续；还需检查收敛的一致性。

	在$[1,+\infty)$上，$\sin x$的原函数有界且$1/x$单调趋于零，故由Dirichlet判别法，$\int_0^\infty(\sin x)/x\d x$收敛（$x=0$为可去奇点）；$\forall x,y\geqslant 0,|\e^{-xy}|\leqslant 1$且$\e^{-xy}$关于$x$递减，由Abel判别准则可知$F(y)$关于$y\in[0,+\infty)$一致收敛，因此$F(y)$在该区间连续，包括端点$y=0$。

	任取$a>0,\forall x\geqslant 0,y\geqslant a,$可知$\abs{-\sin x\e^{-xy}}\leqslant \e^{-xy}\leqslant \e^{-ax}$，且$\int_0^\infty\e^{-ax}\d x=1/a<\infty$，由Weierstrass判别准则知，积分
	\[
		\int\zti\pv {}y\frac{\sin x}x\e^{-xy}\d x=-\int\zti\sin x\e^{-xy}\d x
	\]
	一致收敛，因此$F(y)$在$[a,+\infty)$上可导
	\begin{align*}
        F'(y)&=-\int_0^\infty\sin x\,\e^{-yx}\d x
        =-\Im\int_0^\infty\e^{-(y-\i)x}\d x\\
        &=-\Im\frac1{y-\i}=-\frac1{1+y^2}.
    \end{align*}
	因为$a>0$任意，故$\forall y>0,$
	\[
		F'(y)=-\frac1{1+y^2},\quad F(y)=-\arctan y+C.
	\]
	又$y\to+\infty,$
	\[
		|F(y)|\leqslant\int\zti\frac{|\sin x|}x\e^{-xy}\d x\leqslant\int\zti\e^{-xy}\d x=\frac1y\to0.
	\]
	故$\lim_{y\to+\infty}F(y)=C-\frac{\pi}2=0$，即$C=\frac{\pi}2.$

	又$F(y)$在$[0,+\infty)$上连续，故$\forall y\geqslant 0,F(y)=\frac{\pi}2-\arctan y$。特别地，
	\[
		F(0)=\int\zti\frac{\sin x}x\d x=\frac{\pi}2.
	\]
\end{example}
\newpage
